% ----------------------------------------------------------------------
%  3.5. Электрические свойства; 3.5.1. Основы теории электропроводности
% ----------------------------------------------------------------------
\begin{orientir}
Электросопротивление~— одно из самых чувствительных свойств металла:
по нему следят за полиморфными превращениями, старением, упорядочением,
магнитными переходами, за точечными дефектами после закалки,
деформации и облучения (п.~2.1). Чтобы понимать такие измерения, нужно
знать, от чего зависит сопротивление. Раздел идёт так: модели
проводимости (3.5.1), зависимость от температуры (3.5.2), эффект Холла
(3.5.3), влияние состава, деформации и фазовых превращений
(3.5.4--3.5.6) и сверхпроводимость (3.5.7).
\end{orientir}

\subsection{Основы электронной теории электропроводности}\label{ssec:conduct}

По закону Ома плотность тока пропорциональна напряжённости поля:
\begin{equation}\label{eq:ohm}
\vb j=\sigma\vb E ,\qquad \rho=\frac1\sigma ,\qquad R=\rho\frac lS .
\end{equation}
Удельная проводимость $\sigma$ и удельное сопротивление $\rho$~—
свойства материала, не зависящие от формы образца. У металлов при
комнатной температуре $\rho\approx10^{-8}$--$10^{-7}$~Ом$\cdot$м, т.\,е.
1--10~мкОм$\cdot$см; у меди, одного из лучших проводников,~—
1,55~мкОм$\cdot$см (при 0\,°C).

\paragraph{Модель Друде}
В поле электроны приобретают направленное движение~— \term{дрейф}~— со
средней скоростью $v_\text{д}$. Её ограничивают столкновения с
центрами рассеяния: искажениями решётки, примесями, тепловыми
колебаниями (фононами). Характерное время~— \term{время релаксации}
$\tau$ (время, за которое направленная скорость после выключения поля
спадает в $e$ раз); для чистых металлов $\tau\sim10^{-14}$~с.
\begin{shagi}
\item Поле сообщает электрону ускорение $\d v/\d t=-eE/m$.
\item Если за время $\tau$ электрон «забывает» направленную скорость и
набирает её заново, скорость дрейфа $v_\text{д}=-eE\tau/m$ (против поля:
заряд электрона отрицателен).
\item При концентрации электронов $n$ плотность тока
$j=n(-e)v_\text{д}=ne^2\tau E/m$.
\end{shagi}
\begin{keyeq}\label{eq:drude}
\sigma=\frac{ne^2\tau}{m}.
\end{keyeq}
Друде считал, что электрон сталкивается с ионами, а между
столкновениями ни с чем не взаимодействует; тогда $\tau=l/v$, где $l$~—
длина свободного пробега, и $\sigma=ne^2l/(mv)$. Электроны он считал
классическим газом, у которого средняя квадратичная скорость
$v=\sqrt{3kT/m}$:
\[
\sigma=\frac{ne^2l}{\sqrt{3kTm}} .
\]
При комнатной температуре $v\approx10^7$~см/с, $\tau\sim10^{-14}$~с,
откуда $l\sim1$~нм~— порядка межатомного расстояния\ispr{В оригинале
(с.~489) эти оценки отнесены к 77~К, но приведены $\tau\approx(1$--$20)\cdot10^{-14}$~с
и $l\approx0{,}1$--$1$~нм, что несовместимо: при $v\approx10^7$~см/с такие
$\tau$ дают $l\approx1$--$20$~нм. Значения $l\sim0{,}1$--$1$~нм получаются
при комнатной температуре.}, что вроде бы подтверждает картину
столкновений с ионами. Но модель предсказывает $\rho\propto\sqrt T$, а на
опыте у металлов $\rho\propto T$, и сопротивление растёт гораздо
сильнее.

\paragraph{Квантовая поправка: работают только фермиевские электроны}
Электронный газ вырожден (п.~\ref{ssec:free-el}): все состояния ниже
$\EF$ заняты. Поле может перевести на свободные уровни лишь электроны
вблизи уровня Ферми. Поэтому в электропроводности (как и в
теплоёмкости и теплопроводности) участвуют только они, и движутся они
со скоростью Ферми
\begin{equation}\label{eq:vF}
v_\text{F}=\sqrt{2\EF/m}\approx10^8~\text{см/с},
\end{equation}
а не с тепловой скоростью. Кроме того, электрон не всегда теряет
направленную скорость за одно столкновение: на малоугловых
столкновениях (например, с низкочастотными фононами) для этого нужно
в среднем $\nu$ столкновений. Путь до полной потери направленной
скорости~— \term{транспортная длина свободного пробега} $L=\nu l$:
\begin{keyeq}\label{eq:sigma-quantum}
\sigma=\frac{ne^2L_\text{F}}{mv_\text{F}}=\frac{ne^2\nu l_\text{F}}{mv_\text{F}}.
\end{keyeq}
Формула похожа на классическую, но $v_\text{F}$ от температуры почти не
зависит, а длина пробега у чистых металлов оказывается в десятки и сотни
раз больше межатомного расстояния. Идеальная решётка электроны не
рассеивает (п.~\ref{ssec:bands}); рассеивают только нарушения
периодичности.

\paragraph{Смещение сферы Ферми}
Квантовую картину удобно представить в $\vb k$-пространстве
(рис.~\ref{fig:fermi-shift}). Поле действует на все электроны:
$\hbar\,\d\vb k/\d t=-e\vb E$, и без столкновений сфера Ферми
сдвигалась бы равномерно. Столкновения возвращают электроны назад, и
устанавливается стационарный сдвиг $\delta\vb k=-e\vb E\tau/\hbar$. Ток
создают лишь электроны у «переднего» края смещённой сферы, которым не
нашлось пары с противоположной скоростью; электроны в общей части двух
сфер (внутри) фактически не изменили импульсов. Плотность тока
$j=ne^2\tau E/m=ne^2L_\text{F}E/(mv_\text{F})$~— закон Ома.

\begin{figure}[ht]
\centering
\begin{tikzpicture}[scale=0.85]
  \fill[fmBlue!12] (0,0) circle (1.6);
  \fill[fmRed!25] (0.35,0) circle (1.6);
  \begin{scope}\clip (0,0) circle (1.6); \fill[fmBlue!12] (0.35,0) circle (1.6); \end{scope}
  \draw[fmBlue, thick] (0,0) circle (1.6);
  \draw[fmRed, thick, dashed] (0.35,0) circle (1.6);
  \draw[-{Stealth[length=2mm]}, fmGray] (-2.3,0) -- (2.6,0) node[right, font=\small]{$k_x$};
  \draw[-{Stealth[length=2mm]}, fmGray] (0,-2.1) -- (0,2.2) node[above, font=\small]{$k_y$};
  \draw[-{Stealth[length=2mm]}, thick] (0,-1.9) -- (0.35,-1.9) node[right, font=\small]{$\delta k$};
  \draw[-{Stealth[length=2.4mm]}, line width=1pt, fmGreen!60!black] (2.2,1.6) -- (1.2,1.6) node[left, font=\small]{$\vb E$};
  \node[font=\small, fmBlue, anchor=east] at (-1.65,1.25) {1};
  \node[font=\small, fmRed, anchor=west] at (1.95,1.2) {2};
\end{tikzpicture}
\caption{Смещение сферы Ферми в электрическом поле: 1~— без поля,
2~— в стационарном состоянии с током. Изменились лишь состояния у
поверхности Ферми (розовый серп); электроны в общей части двух сфер
вклада в ток не дают}
\label{fig:fermi-shift}
\end{figure}

\paragraph{Эффективная масса}
Модель свободных электронов объясняет проводимость там, где свободные
электроны есть, но не отвечает, \emph{почему} одни вещества~—
проводники, а другие~— изоляторы. Ответ даёт зонная теория
(п.~\ref{ssec:bands}). Влияние периодического поля учитывают, заменяя
массу электрона \term{эффективной массой}:
\begin{keyeq}\label{eq:meff}
\frac1{m_\text{эфф}}=\frac1{\hbar^2}\frac{\d^2E}{\d k^2},\qquad
\sigma=\frac{ne^2\tau}{m_\text{эфф}} .
\end{keyeq}
Эффективная масса определяется кривизной закона дисперсии $E(k)$
(рис.~\ref{fig:meff}).
\begin{itemize}
\item В центре зоны $E(k)$~— почти парабола свободного электрона,
$m_\text{эфф}\approx m>0$.
\item В точке перегиба $P'$ кривизна равна нулю: $m_\text{эфф}\to\infty$,
поле электрон больше не ускоряет.
\item Ближе к границе зоны кривизна отрицательна, $m_\text{эфф}<0$: поле
\emph{тормозит} электрон (его волна всё сильнее отражается решёткой).
\end{itemize}
Чем сильнее электрон связан с решёткой, тем больше эффективная масса.
У переходных металлов $d$-электроны связаны сильнее $s$-электронов и
имеют большую эффективную массу~— поэтому переходные металлы проводят
ток хуже.

\begin{figure}[ht]
\centering
\begin{tikzpicture}
\begin{groupplot}[group style={group size=1 by 3, vertical sep=4mm, xlabels at=edge bottom, xticklabels at=edge bottom},
  width=0.62\textwidth, height=3.3cm, xmin=0, xmax=1, axis lines=left, ytick={0}, xtick={0.5,1},
  xticklabels={$k_{P'}$,$\pi/a$}, xlabel={$k$}, every axis y label/.style={at={(ticklabel cs:0.5)}, rotate=90, anchor=south}]
  \nextgroupplot[ylabel={$E$}, ymin=0, ymax=2.2]
    \addplot[fmBlue, thick, domain=0:1, samples=80] {1-cos(deg(pi*x))};
    \addplot[fmGray, dashed, domain=0:0.62, samples=40] {(pi*x)^2/2};
    \draw[fmGray, dotted] (axis cs:0.5,0) -- (axis cs:0.5,2.2);
  \nextgroupplot[ylabel={$v$}, ymin=0, ymax=1.15]
    \addplot[fmGreen!60!black, thick, domain=0:1, samples=80] {sin(deg(pi*x))};
    \draw[fmGray, dotted] (axis cs:0.5,0) -- (axis cs:0.5,1.15);
  \nextgroupplot[ylabel={$m_\text{эфф}$}, ymin=-6, ymax=6, ytick={0}]
    \addplot[fmRed, thick, domain=0:0.46, samples=60] {1/cos(deg(pi*x))};
    \addplot[fmRed, thick, domain=0.54:1, samples=60] {1/cos(deg(pi*x))};
    \draw[fmGray] (axis cs:0,0) -- (axis cs:1,0);
    \draw[fmGray, dotted] (axis cs:0.5,-6) -- (axis cs:0.5,6);
\end{groupplot}
\end{tikzpicture}
\caption{Энергия $E$, групповая скорость $v=\hbar^{-1}\d E/\d k$ и
эффективная масса электрона в зоне (модельный закон $E\propto1-\cos ka$;
пунктир~— парабола свободного электрона). В точке перегиба $P'$
скорость максимальна, а эффективная масса обращается в бесконечность и
меняет знак}
\label{fig:meff}
\end{figure}

\paragraph{Проводники, изоляторы, полупроводники}
Если зона заполнена \emph{частично}, состояния с $\vb k$ и $-\vb k$
заселены одинаково, и тока нет; поле смещает распределение
(рис.~\ref{fig:band-field},\,\textit{а, б}), и ток появляется. Если зона
заполнена \emph{целиком} (\textit{в}), смещать нечего: все состояния
заняты, суммарный импульс электронов остаётся нулевым, и тока нет~—
кристалл изолятор. На языке плотности состояний
(рис.~\ref{fig:band-filling}): у проводника зона заполнена частично
или зоны перекрываются; у диэлектрика и полупроводника нижняя зона
заполнена, верхняя пуста, и их разделяет \term{запрещённая зона}
шириной $\Delta E$.

\begin{figure}[ht]
\centering
\begin{tikzpicture}[scale=0.62]
  \foreach \sh/\dx/\full/\lab in {0/0/0/{\textit{а}~— без поля}, 5.2/0.35/0/{\textit{б}~— в поле}, 10.4/0/1/{\textit{в}~— зона заполнена}}{
    \begin{scope}[shift={(\sh,0)}]
      \ifnum\full=1 \fill[fmBlue!30] (-2,-2) rectangle (2,2);
      \else \fill[fmBlue!30] (\dx,0) circle (1.3); \fi
      \draw[atomLine, line width=0.7pt] (-2,-2) rectangle (2,2);
      \draw[fmGray, -{Stealth[length=1.6mm]}] (-2.3,0) -- (2.4,0);
      \draw[fmGray] (0,-2.2) -- (0,2.2);
      \node[font=\small] at (0,-2.75) {\lab};
    \end{scope}}
  \node[font=\scriptsize] at (1.65,1.65) {$\frac\pi a$};
\end{tikzpicture}
\caption{Электроны (синее) в первой зоне Бриллюэна квадратной решётки:
\textit{а}~— частично заполненная зона без поля, токи $\pm\vb k$
компенсируются; \textit{б}~— в поле распределение смещается, течёт ток;
\textit{в}~— в заполненной зоне смещать нечего, тока нет}
\label{fig:band-field}
\end{figure}

\paragraph{Полупроводники}
Они заслуживают отдельного разговора: именно на них построена вся
электроника.
\begin{opredelenie}
\term{Полупроводник}~— твёрдое тело, у которого при $T=0$ высшая
занятая (\term{валентная}) зона полностью заполнена, следующая
(\term{зона проводимости}) пуста, а разделяющая их щель $\Delta E$
невелика.
\end{opredelenie}
Электроны валентной зоны полупроводника легко возбуждаются нагревом
или светом и переходят в зону проводимости, оставляя в валентной зоне
\term{дырки}~— носители положительного заряда. Если носители возникают
за счёт теплового возбуждения через щель~— это \term{собственная}
проводимость; если за счёт примесей~— \term{примесная}.

Число возбуждённых электронов (равное числу дырок) в собственном
полупроводнике
\begin{keyeq}\label{eq:ni}
n_i\approx4{,}8\cdot10^{15}\,T^{3/2}\exp\!\left(-\frac{\Delta E}{2kT}\right)~\text{см}^{-3}
\end{keyeq}
быстро растёт с температурой (рис.~\ref{fig:ni}), а с ним растёт и
проводимость~— в отличие от металлов. Разница между полупроводником и
диэлектриком количественная: у диэлектрика щель так велика, что
заметная проводимость появилась бы лишь при температурах выше
температуры испарения. Условную границу проводят при $\Delta E$ около
2--3~эВ\ispr{В оригинале (с.~495) граница указана как $\Delta E>1$~эВ, но
тогда кремний ($\Delta E=1{,}11$~эВ) пришлось бы считать диэлектриком;
обычно к полупроводникам относят вещества с $\Delta E$ до 2--3~эВ.}.

\begin{primer}[Много ли электронов в зоне проводимости]
При $\Delta E=1$~эВ и $T=293$~К ($kT=0{,}025$~эВ) по \eqref{eq:ni}:
$n_i\approx4{,}8\cdot10^{15}\cdot293^{3/2}\cdot e^{-20}\approx5\cdot10^{10}$~см$^{-3}=5\cdot10^{16}$~м$^{-3}$\isprmark. А атомов в 1~м$^3$ около $5\cdot10^{28}$: в
зону проводимости переходит ничтожная доля~— порядка $10^{-12}$~—
валентных электронов.
\end{primer}
\isprtext{В оригинале (с.~495) приведено $n\sim5\cdot10^{19}$~м$^{-3}$;
расчёт по формуле (3.121) оригинала даёт $5\cdot10^{10}$~см$^{-3}$, т.\,е.
$5\cdot10^{16}$~м$^{-3}$.}

\begin{figure}[ht]
\centering
\begin{tikzpicture}
\begin{axis}[width=0.62\textwidth, height=5.2cm, xmin=200, xmax=700, ymin=1e6, ymax=1e18,
  ymode=log, xlabel={$T$, К}, ylabel={$n_i$, см$^{-3}$}, axis lines=left,
  legend style={font=\scriptsize, at={(0.98,0.05)}, anchor=south east, draw=none, fill=none}]
  \addplot[fmBlue, thick, domain=200:700, samples=60] {4.8e15*x^1.5*exp(-0.66/(2*8.617e-5*x))};
  \addlegendentry{Ge, $\Delta E=0{,}66$~эВ}
  \addplot[fmRed, thick, domain=200:700, samples=60] {4.8e15*x^1.5*exp(-1.11/(2*8.617e-5*x))};
  \addlegendentry{Si, $\Delta E=1{,}11$~эВ}
  \draw[fmGray, dotted] (axis cs:293,1e6) -- (axis cs:293,1e18);
\end{axis}
\end{tikzpicture}
\caption{Концентрация собственных носителей в германии и кремнии
(оценка по \eqref{eq:ni} с массой свободного электрона): при
нагреве от комнатной температуры на 200--300~К она растёт на несколько
порядков}
\label{fig:ni}
\end{figure}

Для сравнения: у высокочистых Si и Ge щель 1,11 и 0,66~эВ, у алмаза~—
около 5,5~эВ\ispr{В оригинале для алмаза указано $\Delta E=6$~эВ; ширина
запрещённой зоны алмаза~— 5,47~эВ.}, это один из лучших диэлектриков.

\paragraph{Примесные полупроводники}
Число носителей резко увеличивают примеси (рис.~\ref{fig:dopants}).
Германию при 20\,°C достаточно одного атома примеси на $10^9$ атомов,
чтобы проводимость стала примесной.
\begin{itemize}
\item \term{Донорные} примеси ($n$-тип). Атом V группы (P, As, Sb) в
решётке германия или кремния (IV группа) приносит на один электрон
больше, чем нужно для четырёх ковалентных связей. Этот электрон связан
слабо; его уровень лежит в щели чуть ниже зоны проводимости~— на
$\Delta E_d\approx0{,}01$~эВ в Ge и $\approx0{,}04$~эВ в Si, и уже при
комнатной температуре он переходит в зону проводимости.
\item \term{Акцепторные} примеси ($p$-тип). Атом III группы (B, Al, Ga)
приносит на один электрон меньше: в связи остаётся незанятое место~—
уровень чуть выше валентной зоны (на $\Delta E_a$). Электроны валентной
зоны легко переходят на него, оставляя в валентной зоне подвижные дырки.
\end{itemize}

\begin{figure}[ht]
\centering
\begin{tikzpicture}[scale=0.8]
  \foreach \sh/\t/\lab in {0/d/{\textit{а}~— донор ($n$-тип)}, 6/a/{\textit{б}~— акцептор ($p$-тип)}}{
    \begin{scope}[shift={(\sh,0)}]
      \fill[fmBlue!35] (0,0) rectangle (4,1.2);
      \fill[fmGray!12] (0,2.8) rectangle (4,4);
      \draw[atomLine] (0,0) rectangle (4,1.2) (0,2.8) rectangle (4,4);
      \node[font=\scriptsize] at (2,0.6) {валентная зона};
      \node[font=\scriptsize] at (2,3.4) {зона проводимости};
      \node[font=\small] at (2,-0.5) {\lab};
    \end{scope}}
  \foreach \x in {0.5,1.3,2.1,2.9,3.7}{\draw[fmRed, line width=1.4pt] (\x-0.25,2.55) -- (\x+0.25,2.55);}
  \draw[dimline] (4.25,2.55) -- node[right, font=\scriptsize]{$\Delta E_d$} (4.25,2.8);
  \foreach \x in {6.5,7.3,8.1,8.9,9.7}{\draw[fmGreen!60!black, line width=1.4pt] (\x-0.25,1.45) -- (\x+0.25,1.45);}
  \draw[dimline] (10.25,1.2) -- node[right, font=\scriptsize]{$\Delta E_a$} (10.25,1.45);
  \draw[-{Stealth[length=1.6mm]}, fmRed] (1.3,2.62) -- (1.3,3.05);
  \fill[fmRed] (1.3,3.15) circle (0.07);
  \draw[-{Stealth[length=1.6mm]}, fmGreen!60!black] (8.1,1.0) -- (8.1,1.38);
  \draw[fmGreen!60!black] (8.1,0.85) circle (0.08);
\end{tikzpicture}
\caption{Примесные уровни в запрещённой зоне: \textit{а}~— донорные
уровни под зоной проводимости, с них электроны легко переходят в зону;
\textit{б}~— акцепторные уровни над валентной зоной, на них переходят
электроны, оставляя в валентной зоне дырки (кружок)}
\label{fig:dopants}
\end{figure}
